% Generated draft: requires mathematical and source review.
\documentclass[11pt]{article}
\usepackage[margin=25mm]{geometry}
\usepackage{fontspec}
\setmainfont{Latin Modern Roman}
\usepackage{amsmath,amssymb,mathtools,unicode-math}
\setmathfont{Latin Modern Math}
\usepackage{enumitem,needspace,xurl,hyperref}
\setlength{\parindent}{0pt}
\setlength{\parskip}{5pt}
\setlength{\emergencystretch}{3em}
\newcommand{\N}{\mathbb{N}}
\newcommand{\Z}{\mathbb{Z}}
\newcommand{\Q}{\mathbb{Q}}
\newcommand{\R}{\mathbb{R}}
\newcommand{\C}{\mathbb{C}}
\newcommand{\F}{\mathbb{F}}
\newcommand{\A}{\mathbb{A}}
\newcommand{\PP}{\mathbb{P}}
\newcommand{\E}{\mathbb{E}}
\newcommand{\eps}{\varepsilon}
\newcommand{\dd}{\,\mathrm{d}}
\begin{document}
\section*{UnsolvedMath ID: 30004296}
\textbf{Source problem code:} catalogue entry
\subsection{Definitions and mathematical statement}
The supplied problem concerns a choice of function
\[
f(n)\in\{\sqrt n,\ n,\ n^{3/2},\ n^2,\ n^3\},
\]
with one function chosen and applied to every entry of an integer matrix \(m=(m_{ij})\). The matrix dimensions, admissible entry values, and any size cutoff are not specified in the record.

Conditionally on choosing positive integers \(r,s\) and an admissible set \(\mathcal A\subseteq\Z^{r\times s}\) on which the chosen function is real-valued, put
\[
M_f(m)=(f(m_{ij}))_{1\le i\le r,\,1\le j\le s}.
\]
For \(1\le k\le\min(r,s)\), define
\[
\Delta_k(M_f(m))=
\max_{\substack{I\subseteq\{1,\ldots,r\},\ J\subseteq\{1,\ldots,s\}\\ |I|=|J|=k}}
\left|\det\bigl(f(m_{ij})\bigr)_{i\in I,\,j\in J}\right|,
\]
where row and column indices are taken in increasing order.

The requested counting function counts tuples satisfying prescribed conditions \(\Delta_k(M_f(m))\asymp D_k\). More explicitly, once a set of constrained orders \(K\subseteq\{1,\ldots,\min(r,s)\}\), positive scales \(D_k\), and comparison constants \(0<c_k\le C_k<\infty\) have been specified, this would mean
\[
\mathcal N_{f,\mathcal A,K}(\mathbf D;\mathbf c,\mathbf C)
=
\#\left\{m\in\mathcal A:
 c_kD_k\le\Delta_k(M_f(m))\le C_kD_k
 \text{ for every }k\in K\right\}.
\]
The source asks what upper bounds can be obtained for this counting function. The parameters above expose missing specifications; they are not additional hypotheses attributed to the source.

No definite quantitative bound is requested in the supplied text. In particular, the admissible set, the constrained minor orders, the meaning and permitted parameter dependence of the comparison constants, the asymptotic regime, and treatment of zero scales remain unresolved. Without an entry restriction the count need not be finite. For the square-root and \(n^{3/2}\) choices, a real-valued interpretation also requires an admissible domain, which the source does not state.
\subsection{Short English statement}
Choose one of the five listed functions, apply it entrywise to an integer matrix, and count matrices whose largest absolute minors have prescribed sizes. The problem asks for upper bounds, but the supplied record omits the matrix dimensions, entry range, and precise scale conventions needed to specify a finite counting problem and its intended bounds.
\subsection{Sources}
\begin{itemize}
\item Supplied local UnsolvedMath registry record, permanent ID \(30004296\), title: \emph{Counting Integer Matrices with Prescribed Minor Scales}; exact source problem code: \texttt{catalogue entry}. Collection: \emph{Oberwolfach Reports Open Problems}. No individual proposer is supplied. The proposed year is recorded as 2019.
\item Supplied source citation: \emph{Analytic Number Theory (2020)}, Oberwolfach Reports, pp.\ 3141--3205, DOI \texttt{10.4171/owr/2019/50}. Source URL and supplied literature reference labelled ``original OWR report'': \url{https://doi.org/10.4171/owr/2019/50}. The record describes these reports as documenting mathematical workshops and their open-problem discussions; no specific problem page or individual attribution is supplied.
\item Literature assessment dated exactly \texttt{2026-08-22}: the supplied curation labels the problem open, describes it as restriction-theoretic, and states that existing determinant and restriction estimates do not provide the requested general upper bounds. This is supplied dated curation, not an independently verified assessment. The record explicitly requests independent verification of time-sensitive status and cited claims.
\item The supplied statement status is \texttt{corrected\_verified}. Its accompanying qualification says that only formula delimiters, notation, grammar, and typography were repaired, with mathematical content and scope preserved. This label does not establish mathematical completeness or correctness of this draft.
\end{itemize}
\end{document}
